\documentclass[12pt,a4paper]{article}

% --- Настройка языка и шрифтов для LuaLaTeX ---
\usepackage{fontspec}
\defaultfontfeatures{Ligatures=TeX,Scale=MatchLowercase}

\setmainfont{Alice-Regular}[
    Path=font/,
    Extension=.ttf,
    BoldFont = Alice-Regular,
    ItalicFont = Alice-Regular,
    BoldItalicFont = Alice-Regular,
    BoldFeatures = {FakeBold=1.6},
    ItalicFeatures = {FakeSlant=0.2},
    BoldItalicFeatures = {FakeBold=1.6,FakeSlant=0.2}
]

\newfontfamily\cyrillicfont{Alice-Regular}[
    Path=font/,
    Extension=.ttf,
    BoldFont = Alice-Regular,
    ItalicFont = Alice-Regular,
    BoldItalicFont = Alice-Regular,
    BoldFeatures = {FakeBold=1.6},
    ItalicFeatures = {FakeSlant=0.2},
    BoldItalicFeatures = {FakeBold=1.6,FakeSlant=0.2}
]

\usepackage{polyglossia}
\setmainlanguage{russian}
\setsansfont{TeX Gyre Heros}
\setmonofont{TeX Gyre Cursor}

\usepackage{graphicx}        					% Графика
\usepackage{amsmath, amsfonts, amssymb, amsthm}	% Математика
\usepackage{mathtools}
\usepackage{chemfig}
\usepackage[version=4,arrows=pgf-filled,
            textfontname=sffamily,
            mathfontname=mathsf]{mhchem}
\usepackage[protrusion=false]{microtype}       % Микротипографика
\usepackage{braket} 						% Braket нотация для квантовой механики
\usepackage{unicode-math} 					% Современный пакет для математики в LuaLaTeX
\usepackage{longtable}       				% Таблицы с разрывом страниц
\usepackage{tikz}           				% Рисование
\usepackage[europeanresistors]{circuitikz}	% Рисование электрических схем в европейском стиле
\usepackage{listings}						% Для отображения кода
\usepackage{tabularx}       				% Для таблиц во всю ширину
\usepackage{array}		  					% Для расширенных возможностей таблиц
\usepackage{float}          				% Для точного позиционирования фигур
\usepackage{pgfplots}                          % Для построения графиков
\usepackage{caption}                           % Для настройки подписей к рисункам и таблицам
\usepackage{xfp}                               % Для вычислений в LaTeX
\usepackage{luacode}                           % Для выполнения Lua-кода в документе    
\pgfplotsset{compat=1.18}
\usepackage[paper=letterpaper,
            top=20mm,
            bottom=20mm,
            left=10mm,
            right=10mm,
            footskip=0.4in 
]{geometry}

\usetikzlibrary{arrows.meta}

% --- Колонтитулы ---
\usepackage{fancyhdr}
\pagestyle{fancy}
\fancyhf{}

% --- Теоремы и окружения ---
\theoremstyle{plain}
\newtheorem{theorem}{Theorem}[section]
\newtheorem{lemma}[theorem]{Lemma}
\newtheorem{proposition}[theorem]{Proposition}
\newtheorem{corollary}[theorem]{Corollary}

\theoremstyle{definition}
\newtheorem{definition}[theorem]{Definition}

\theoremstyle{remark}
\newtheorem{remark}[theorem]{Remark}

\numberwithin{equation}{section}


\renewcommand{\headrulewidth}{0.4pt}
\renewcommand{\footrulewidth}{0.4pt}
\renewcommand\lstlistingname{Исходный код}

\lstdefinestyle{main}{
	backgroundcolor=\color{gray!8},   
	commentstyle=\color{green!50!black},
	keywordstyle=\color{blue!80!black},
	numberstyle=\small\color{darkgray},
	stringstyle=\rmfamily\color{orange!90!black},
	basicstyle=\ttfamily\small,
	breakatwhitespace=false,         
	breaklines=true,                   
	keepspaces=true,                 
	numbers=left,                    
	numbersep=8pt,                  
	showspaces=false,                
	showstringspaces=false,
	showtabs=false,                  
	tabsize=4
}	

\lstset{style=main}

\newcolumntype{C}{>{\centering\arraybackslash}X}
\renewcommand{\arraystretch}{1.25}

\begin{luacode*}
    data = require("data")
    include = require("include")
    handle = require("handle")
\end{luacode*}

\begin{document}

\section{\large\textnormal{Для резкого $p$-$n$-перехода вычислить \\
- контактную разность потенциалов $\varphi_k$, \\
- полную ширину ОПЗ $l_0$, \\
- соотношение ширин ОПЗ в $p$ и $n$-областях $\frac{l_{p_0}}{l_{n_0}}$, \\
- максимальное электрическое поле в ОПЗ $E_{m_0}$ \\
при заданных в соответствии с вариантом задания $N_\text{э}$ и $N_\text{б}$.}}

\begin{figure}[H]
    \centering
    \begin{minipage}[t]{0.3\textwidth}
    \vspace{0em}
    \begin{tabularx}{\textwidth}{|X|}
        \hline
        \multicolumn{1}{|C|}{Дано} \\
        \hline
        $N_\text{э} = 10^{20}$ см$^{-3}$ \\ 
        $N_\text{б} = 10^{18}$ см$^{-3}$ \\
        $n_i = 1.45 \cdot 10^{10}$ см$^{-3}$ \\
        $\varepsilon_0 = 8.854 \cdot 10^{-14}$ Ф/см \\
        $\varepsilon_\text{Si} = 11.7$ \\
        $T = 300$ К \\
        \hline
        \multicolumn{1}{|C|}{Найти} \\
        \hline
        $\varphi_k$ — ? \\
         $l_0$ — ? \\
         $\frac{l_{p_0}}{l_{n_0}}$ — ? \\
         $E_{m_0}$ — ? \\
        \hline
    \end{tabularx}

    \vspace{1em}
    Контактная разность потенциалов $\varphi_k$ определяется как
    \begin{equation*}
        \varphi_k = \frac{kT}{e}\ln\left(\frac{N_\text{э}N_\text{б}}{n_i^2}\right)
    \end{equation*}
    
    Ширина ОПЗ $l_0$ определяется как
    \begin{equation*}
        l_0 = \sqrt{\frac{2\varepsilon\varepsilon_0\varphi_k}{e}\left(\frac{1}{N_\text{э}} + \frac{1}{N_\text{б}}\right)}  
    \end{equation*}

    Соотношение ширин ОПЗ в $p$ и $n$-областях $\frac{l_{p_0}}{l_{n_0}}$ определяется как
    \begin{equation*}
        \frac{l_{p_0}}{l_{n_0}} = \frac{N_\text{б}}{N_\text{э}}
    \end{equation*}

    Максимальное электрическое поле в ОПЗ $E_{m_0}$ определяется как
    \begin{equation*}
        E_{m_0} = \frac{2\varphi_k}{l_0}
    \end{equation*}

    \end{minipage}\hfill\begin{minipage}[t]{0.7\textwidth}
        \vspace{0em}
        \begin{tabularx}{\textwidth}{C|}
            \hline
            Решение \\
            \hline
        \end{tabularx}

        \vspace{0.5em}
        \hspace{0.025\textwidth}
        \begin{minipage}[t]{0.95\textwidth}
            Считаем 
            \begin{equation*}
                \begin{aligned}
                    \varphi_k &= \frac{kT}{e}\ln\left(\frac{N_\text{э}N_\text{б}}{n_i^2}\right) = \\
                    &= \directlua{
                        tex.sprint("\\frac{"
                                   .. include.value_with_mantice(data.k, 2, 0, 0) .. 
                                   " \\cdot "
                                   .. include.value_with_mantice(data.T, 2, 0, 0) .. 
                                   "}{"
                                   .. include.value_with_mantice(data.e, 2, 0, 0) .. 
                                   "} \\ln\\left(\\frac{"
                                   .. include.value_with_mantice(data.N_e, 2, 0, 0) .. 
                                   " \\cdot "
                                   .. include.value_with_mantice(data.N_b, 2, 0, 0) .. 
                                   "}{("
                                   .. include.value_with_mantice(data.n_i, 2, 0, 0) .. 
                                   ")^2}\\right) = "
                                   .. include.value_with_mantice(data.phi_k, 2, 0, 0))
                    } \text{ В}
                \end{aligned}
            \end{equation*}
            \begin{equation*}
                \begin{aligned}
                    l_0 &= \sqrt{\frac{2\varepsilon\varepsilon_0\varphi_k}{e}\left(\frac{1}{N_\text{э}} + \frac{1}{N_\text{б}}\right)} = \\ 
                    &= \directlua{
                        tex.sprint("\\sqrt{\\frac{2 \\cdot "
                                   .. include.value_with_mantice(data.epsilon_SI, 2, 0, 0) .. 
                                   " \\cdot "
                                   .. include.value_with_mantice(data.epsilon_0, 2, 0, 0) .. 
                                   " \\cdot "
                                   .. include.value_with_mantice(data.phi_k, 2, 0, 0) .. 
                                   "}{"
                                   .. include.value_with_mantice(data.e, 2, 0, 0) .. 
                                   "} \\left(\\frac{1}{"
                                   .. include.value_with_mantice(data.N_e, 2, 0, 0) .. 
                                   "} + \\frac{1}{"
                                   .. include.value_with_mantice(data.N_b, 2, 0, 0) .. 
                                   "}\\right)}")
                    } = \\
                    &= \directlua{
                        tex.sprint(include.value_with_mantice(data.l_0, 2, 0, 0))
                    } \text{ см}
                \end{aligned}
            \end{equation*}
            \begin{equation*}
                \frac{l_{p_0}}{l_{n_0}} = \frac{N_\text{б}}{N_\text{э}} = \directlua{
                    tex.sprint("\\frac{"
                               .. include.value_with_mantice(data.N_b, 2, 0, 0) .. 
                               "}{"
                               .. include.value_with_mantice(data.N_e, 2, 0, 0) .. 
                               "} = "
                               .. include.value_with_mantice(data.N_b / data.N_e, 2, 0, 0))
                }
            \end{equation*}
            \begin{equation*}
                E_{m_0} = \frac{2\varphi_k}{l_0} = \directlua{
                    tex.sprint("\\frac{2 \\cdot "
                               .. include.value_with_mantice(data.phi_k, 2, 0, 0) .. 
                               "}{"
                               .. include.value_with_mantice(data.l_0, 2, 0, 0) .. 
                               "} = "
                               .. include.value_with_mantice(data.E_m_0, 2, 0, 0))
                } \text{ В/см}
            \end{equation*}

            Ответ: $\varphi_k = \directlua{
                tex.sprint(include.value_with_mantice(data.phi_k, 2, 0, 0))
            } \text{ В}$, $l_0 = \directlua{
                tex.sprint(include.value_with_mantice(data.l_0, 2, 0, 0))
            } \text{ см}$, $\frac{l_{p_0}}{l_{n_0}} = \directlua{
                tex.sprint(include.value_with_mantice(data.N_b / data.N_e, 2, 0, 0))
            }$, $E_{m_0} = \directlua{
                tex.sprint(include.value_with_mantice(data.E_m_0, 2, 0, 0))
            } \text{ В/см}$.
    \end{minipage}\end{minipage}
\end{figure}

\pagebreak

\section{\large\textnormal{Для линейного $p$-$n$-перехода вычислить \\
- контактную разность потенциалов $\varphi_k$, \\
- полную ширину ОПЗ $l_0$ \\
при заданном в соответствии с вариантом задания градиенте концентрации $N'$.}}

\begin{figure}[H]
    \centering
    \begin{minipage}[t]{0.3\textwidth}
    \vspace{0em}
    \begin{tabularx}{\textwidth}{|X|}
        \hline
        \multicolumn{1}{|C|}{Дано} \\
        \hline
        $N' = 10^{24}$ см$^{-4}$ \\
        $n_i = 1.45 \cdot 10^{10}$ см$^{-3}$ \\
        $\varepsilon_0 = 8.854 \cdot 10^{-14}$ Ф/см \\
        $\varepsilon_\text{Si} = 11.7$ \\
        $T = 300$ К \\
        \hline
        \multicolumn{1}{|C|}{Найти} \\
        \hline
        $\varphi_k$ — ? \\
        $l_0$ — ? \\
        \hline
    \end{tabularx}

    \vspace{1em}
    Контактная разность потенциалов $\varphi_k$ определяется как
    \begin{equation*}
        \varphi_k = \frac{2kT}{e}\ln\left(\frac{N' l_0}{2n_i}\right)
    \end{equation*}

    Ширина ОПЗ $l_0$ определяется как
    \begin{equation*}
        l_0 = \sqrt[3]{\frac{12\varepsilon\varepsilon_0\varphi_k}{e N'}}
    \end{equation*}

    \end{minipage}\hfill\begin{minipage}[t]{0.7\textwidth}
        \vspace{0em}
        \begin{tabularx}{\textwidth}{C|}
            \hline
            Решение \\
            \hline
        \end{tabularx}

        \vspace{0.5em}
        \hspace{0.025\textwidth}
        \begin{minipage}[t]{0.95\textwidth}
            Поскольку $\varphi_k$ и $l_0$ взаимно связаны, расчет проводится методом последовательных приближений. Зададим начальное приближение $\varphi_k^{(0)} = 0.7\text{ В}$.

            Первое приближение:
            \begin{equation*}
                \begin{aligned}
                    l_0^{(0)} &= \sqrt[3]{\frac{12\varepsilon\varepsilon_0\varphi_k^{(0)}}{e N'}} = \\
                    &= \directlua{
                        tex.sprint("\\sqrt[3]{\\frac{12 \\cdot "
                                   .. include.value_with_mantice(data.epsilon_SI, 2, 0, 0) .. 
                                   " \\cdot "
                                   .. include.value_with_mantice(data.epsilon_0, 2, 0, 0) .. 
                                   " \\cdot "
                                   .. include.value_with_mantice(data.phi_k_0_lin, 2, 0, 0) .. 
                                   "}{"
                                   .. include.value_with_mantice(data.e, 2, 0, 0) .. 
                                   " \\cdot "
                                   .. include.value_with_mantice(data.N_prime, 2, 0, 0) .. 
                                   "}} = "
                                   .. include.value_with_mantice(data.l_0_0_lin, 2, 0, 0))
                    } \text{ см}
                \end{aligned}
            \end{equation*}
            \begin{equation*}
                \begin{aligned}
                    \varphi_k^{(1)} &= \frac{2kT}{e}\ln\left(\frac{N' l_0^{(0)}}{2n_i}\right) = \\
                    &= \directlua{
                        tex.sprint("\\frac{2 \\cdot "
                                   .. include.value_with_mantice(data.k, 2, 0, 0) .. 
                                   " \\cdot "
                                   .. include.value_with_mantice(data.T, 2, 0, 0) .. 
                                   "}{"
                                   .. include.value_with_mantice(data.e, 2, 0, 0) .. 
                                   "} \\ln\\left(\\frac{"
                                   .. include.value_with_mantice(data.N_prime, 2, 0, 0) .. 
                                   " \\cdot "
                                   .. include.value_with_mantice(data.l_0_0_lin, 2, 0, 0) .. 
                                   "}{2 \\cdot "
                                   .. include.value_with_mantice(data.n_i, 2, 0, 0) .. 
                                   "}\\right) = "
                                   .. include.value_with_mantice(data.phi_k_1_lin, 2, 0, 0))
                    } \text{ В}
                \end{aligned}
            \end{equation*}

            Второе приближение:
            \begin{equation*}
                \begin{aligned}
                    l_0^{(1)} &= \sqrt[3]{\frac{12\varepsilon\varepsilon_0\varphi_k^{(1)}}{e N'}} = \\
                    &= \directlua{
                        tex.sprint("\\sqrt[3]{\\frac{12 \\cdot "
                                   .. include.value_with_mantice(data.epsilon_SI, 2, 0, 0) .. 
                                   " \\cdot "
                                   .. include.value_with_mantice(data.epsilon_0, 2, 0, 0) .. 
                                   " \\cdot "
                                   .. include.value_with_mantice(data.phi_k_1_lin, 2, 0, 0) .. 
                                   "}{"
                                   .. include.value_with_mantice(data.e, 2, 0, 0) .. 
                                   " \\cdot "
                                   .. include.value_with_mantice(data.N_prime, 2, 0, 0) .. 
                                   "}} = "
                                   .. include.value_with_mantice(data.l_0_1_lin, 2, 0, 0))
                    } \text{ см}
                \end{aligned}
            \end{equation*}
            \begin{equation*}
                \begin{aligned}
                    \varphi_k^{(2)} &= \frac{2kT}{e}\ln\left(\frac{N' l_0^{(1)}}{2n_i}\right) = \\
                    &= \directlua{
                        tex.sprint("\\frac{2 \\cdot "
                                   .. include.value_with_mantice(data.k, 2, 0, 0) .. 
                                   " \\cdot "
                                   .. include.value_with_mantice(data.T, 2, 0, 0) .. 
                                   "}{"
                                   .. include.value_with_mantice(data.e, 2, 0, 0) .. 
                                   "} \\ln\\left(\\frac{"
                                   .. include.value_with_mantice(data.N_prime, 2, 0, 0) .. 
                                   " \\cdot "
                                   .. include.value_with_mantice(data.l_0_1_lin, 2, 0, 0) .. 
                                   "}{2 \\cdot "
                                   .. include.value_with_mantice(data.n_i, 2, 0, 0) .. 
                                   "}\\right) = "
                                   .. include.value_with_mantice(data.phi_k_2_lin, 2, 0, 0))
                    } \text{ В}
                \end{aligned}
            \end{equation*}

            Значения сходятся к:
            \begin{equation*}
                \begin{aligned}
                    l_0 &= \directlua{
                        tex.sprint(include.value_with_mantice(data.l_0_lin, 2, 0, 0))
                    } \text{ см} \\
                    \varphi_k &= \directlua{
                        tex.sprint(include.value_with_mantice(data.phi_k_lin, 2, 0, 0))
                    } \text{ В}
                \end{aligned}
            \end{equation*}

            Ответ: $\varphi_k = \directlua{
                tex.sprint(include.value_with_mantice(data.phi_k_lin, 2, 0, 0))
            } \text{ В}$, $l_0 = \directlua{
                tex.sprint(include.value_with_mantice(data.l_0_lin, 2, 0, 0))
            } \text{ см}$.
    \end{minipage}\end{minipage}
\end{figure}

\pagebreak

\section{\large\textnormal{По результатам предыдущих задач рассчитать значения $C_0$ для ступенчатого и линейного $p$-$n$-переходов площадью $S_1 = 10\text{ см}^2$.}}

\begin{figure}[H]
    \centering
    \begin{minipage}[t]{0.3\textwidth}
    \vspace{0em}
    \begin{tabularx}{\textwidth}{|X|}
        \hline
        \multicolumn{1}{|C|}{Дано} \\
        \hline
        $S_1 = 10$ см$^2$ \\
        $l_{0,\text{резк}} = \directlua{tex.sprint(include.value_with_mantice(data.l_0, 2, 0, 0))}$ см \\
        $l_{0,\text{лин}} = \directlua{tex.sprint(include.value_with_mantice(data.l_0_lin, 2, 0, 0))}$ см \\
        $\varepsilon_0 = 8.854 \cdot 10^{-14}$ Ф/см \\
        $\varepsilon_\text{Si} = 11.7$ \\
        \hline
        \multicolumn{1}{|C|}{Найти} \\
        \hline
        $C_{0,\text{резк}}$ — ? \\
        $C_{0,\text{лин}}$ — ? \\
        \hline
    \end{tabularx}

    \vspace{1em}
    Барьерная емкость $C_0$ при нулевом смещении определяется как
    \begin{equation*}
        C_0 = \frac{\varepsilon\varepsilon_0 S_1}{l_0}
    \end{equation*}

    Для ступенчатого (резкого) перехода:
    \begin{equation*}
        C_{0,\text{резк}} = \frac{\varepsilon\varepsilon_0 S_1}{l_{0,\text{резк}}}
    \end{equation*}

    Для линейного перехода:
    \begin{equation*}
        C_{0,\text{лин}} = \frac{\varepsilon\varepsilon_0 S_1}{l_{0,\text{лин}}}
    \end{equation*}

    \end{minipage}\hfill\begin{minipage}[t]{0.7\textwidth}
        \vspace{0em}
        \begin{tabularx}{\textwidth}{C|}
            \hline
            Решение \\
            \hline
        \end{tabularx}

        \vspace{0.5em}
        \hspace{0.025\textwidth}
        \begin{minipage}[t]{0.95\textwidth}
            Считаем барьерную емкость для ступенчатого (резкого) перехода:
            \begin{equation*}
                \begin{aligned}
                    C_{0,\text{резк}} &= \frac{\varepsilon\varepsilon_0 S_1}{l_{0,\text{резк}}} = \\
                    &= \directlua{
                        tex.sprint("\\frac{"
                                   .. include.value_with_mantice(data.epsilon_SI, 2, 0, 0) .. 
                                   " \\cdot "
                                   .. include.value_with_mantice(data.epsilon_0, 2, 0, 0) .. 
                                   " \\cdot "
                                   .. include.value_with_mantice(data.S_1, 2, 0, 0) .. 
                                   "}{"
                                   .. include.value_with_mantice(data.l_0, 2, 0, 0) .. 
                                   "} = "
                                   .. include.value_with_mantice(data.C_0_abrupt, 2, 0, 0))
                    } \text{ Ф}
                \end{aligned}
            \end{equation*}

            Считаем барьерную емкость для линейного перехода:
            \begin{equation*}
                \begin{aligned}
                    C_{0,\text{лин}} &= \frac{\varepsilon\varepsilon_0 S_1}{l_{0,\text{лин}}} = \\
                    &= \directlua{
                        tex.sprint("\\frac{"
                                   .. include.value_with_mantice(data.epsilon_SI, 2, 0, 0) .. 
                                   " \\cdot "
                                   .. include.value_with_mantice(data.epsilon_0, 2, 0, 0) .. 
                                   " \\cdot "
                                   .. include.value_with_mantice(data.S_1, 2, 0, 0) .. 
                                   "}{"
                                   .. include.value_with_mantice(data.l_0_lin, 2, 0, 0) .. 
                                   "} = "
                                   .. include.value_with_mantice(data.C_0_lin, 2, 0, 0))
                    } \text{ Ф}
                \end{aligned}
            \end{equation*}

            Ответ: $C_{0,\text{резк}} = \directlua{
                tex.sprint(include.value_with_mantice(data.C_0_abrupt, 2, 0, 0))
            } \text{ Ф}$, $C_{0,\text{лин}} = \directlua{
                tex.sprint(include.value_with_mantice(data.C_0_lin, 2, 0, 0))
            } \text{ Ф}$.
    \end{minipage}\end{minipage}
\end{figure}

\pagebreak

\section{\large\textnormal{Для полученных в предыдущих заданиях результатов расчетов параметров резкого $p$-$n$-перехода рассчитать $C(U)$ при $U = -5\text{ В}$ и $U = +0.5\text{ В}$.}}

\begin{figure}[H]
    \centering
    \begin{minipage}[t]{0.3\textwidth}
    \vspace{0em}
    \begin{tabularx}{\textwidth}{|X|}
        \hline
        \multicolumn{1}{|C|}{Дано} \\
        \hline
        $C_{0,\text{резк}} = \directlua{tex.sprint(include.value_with_mantice(data.C_0_abrupt, 2, 0, 0))}$ Ф \\
        $\varphi_k = \directlua{tex.sprint(include.value_with_mantice(data.phi_k, 2, 0, 0))}$ В \\
        $U_1 = -5$ В \\
        $U_2 = +0.5$ В \\
        \hline
        \multicolumn{1}{|C|}{Найти} \\
        \hline
        $C(U_1)$ — ? \\
        $C(U_2)$ — ? \\
        \hline
    \end{tabularx}

    \vspace{1em}
    Зависимость барьерной емкости резкого перехода от напряжения $U$ определяется как
    \begin{equation*}
        C(U) = \frac{C_0}{\sqrt{1 - \frac{U}{\varphi_k}}}
    \end{equation*}
    При обратном смещении ($U < 0$) ширина ОПЗ возрастает, а емкость уменьшается. При прямом смещении ($U > 0$) ширина ОПЗ уменьшается, а емкость возрастает.

    \end{minipage}\hfill\begin{minipage}[t]{0.7\textwidth}
        \vspace{0em}
        \begin{tabularx}{\textwidth}{C|}
            \hline
            Решение \\
            \hline
        \end{tabularx}

        \vspace{0.5em}
        \hspace{0.025\textwidth}
        \begin{minipage}[t]{0.95\textwidth}
            Считаем при обратном смещении $U_1 = -5\text{ В}$:
            \begin{equation*}
                \begin{aligned}
                    C(U_1) &= \frac{C_{0,\text{резк}}}{\sqrt{1 - \frac{U_1}{\varphi_k}}} = \\
                    &= \directlua{
                        tex.sprint("\\frac{"
                                   .. include.value_with_mantice(data.C_0_abrupt, 2, 0, 0) .. 
                                   "}{\\sqrt{1 - \\frac{("
                                   .. include.value_with_mantice(data.U_1, 2, 0, 0) .. 
                                   ")}{"
                                   .. include.value_with_mantice(data.phi_k, 2, 0, 0) .. 
                                   "}}} = "
                                   .. include.value_with_mantice(data.C_abrupt_U1, 2, 0, 0))
                    } \text{ Ф}
                \end{aligned}
            \end{equation*}

            Считаем при прямом смещении $U_2 = +0.5\text{ В}$:
            \begin{equation*}
                \begin{aligned}
                    C(U_2) &= \frac{C_{0,\text{резк}}}{\sqrt{1 - \frac{U_2}{\varphi_k}}} = \\
                    &= \directlua{
                        tex.sprint("\\frac{"
                                   .. include.value_with_mantice(data.C_0_abrupt, 2, 0, 0) .. 
                                   "}{\\sqrt{1 - \\frac{"
                                   .. include.value_with_mantice(data.U_2, 2, 0, 0) .. 
                                   "}{"
                                   .. include.value_with_mantice(data.phi_k, 2, 0, 0) .. 
                                   "}}} = "
                                   .. include.value_with_mantice(data.C_abrupt_U2, 2, 0, 0))
                    } \text{ Ф}
                \end{aligned}
            \end{equation*}

            Ответ: $C(-5\text{ В}) = \directlua{
                tex.sprint(include.value_with_mantice(data.C_abrupt_U1, 2, 0, 0))
            } \text{ Ф}$, $C(+0.5\text{ В}) = \directlua{
                tex.sprint(include.value_with_mantice(data.C_abrupt_U2, 2, 0, 0))
            } \text{ Ф}$.
    \end{minipage}\end{minipage}
\end{figure}

\pagebreak

\section{\large\textnormal{Для полученных в предыдущих заданиях результатов расчетов параметров линейного $p$-$n$-перехода рассчитать $C(U)$ при $U = -5\text{ В}$ и $U = +0.5\text{ В}$.}}

\begin{figure}[H]
    \centering
    \begin{minipage}[t]{0.3\textwidth}
    \vspace{0em}
    \begin{tabularx}{\textwidth}{|X|}
        \hline
        \multicolumn{1}{|C|}{Дано} \\
        \hline
        $C_{0,\text{лин}} = \directlua{tex.sprint(include.value_with_mantice(data.C_0_lin, 2, 0, 0))}$ Ф \\
        $\varphi_k = \directlua{tex.sprint(include.value_with_mantice(data.phi_k_lin, 2, 0, 0))}$ В \\
        $U_1 = -5$ В \\
        $U_2 = +0.5$ В \\
        \hline
        \multicolumn{1}{|C|}{Найти} \\
        \hline
        $C(U_1)$ — ? \\
        $C(U_2)$ — ? \\
        \hline
    \end{tabularx}

    \vspace{1em}
    Зависимость барьерной емкости линейного перехода от напряжения $U$ определяется как
    \begin{equation*}
        C(U) = \frac{C_0}{\sqrt[3]{1 - \frac{U}{\varphi_k}}}
    \end{equation*}
    где показатель степени равен $1/3$, так как для линейного перехода ширина ОПЗ зависит от напряжения как:
    \begin{equation*}
        l(U) \sim \sqrt[3]{\varphi_k - U}
    \end{equation*}

    \end{minipage}\hfill\begin{minipage}[t]{0.7\textwidth}
        \vspace{0em}
        \begin{tabularx}{\textwidth}{C|}
            \hline
            Решение \\
            \hline
        \end{tabularx}

        \vspace{0.5em}
        \hspace{0.025\textwidth}
        \begin{minipage}[t]{0.95\textwidth}
            Считаем при обратном смещении $U_1 = -5\text{ В}$:
            \begin{equation*}
                \begin{aligned}
                    C(U_1) &= \frac{C_{0,\text{лин}}}{\sqrt[3]{1 - \frac{U_1}{\varphi_k}}} = \\
                    &= \directlua{
                        tex.sprint("\\frac{"
                                   .. include.value_with_mantice(data.C_0_lin, 2, 0, 0) .. 
                                   "}{\\sqrt[3]{1 - \\frac{("
                                   .. include.value_with_mantice(data.U_1, 2, 0, 0) .. 
                                   ")}{"
                                   .. include.value_with_mantice(data.phi_k_lin, 2, 0, 0) .. 
                                   "}}} = "
                                   .. include.value_with_mantice(data.C_lin_U1, 2, 0, 0))
                    } \text{ Ф}
                \end{aligned}
            \end{equation*}

            Считаем при прямом смещении $U_2 = +0.5\text{ В}$:
            \begin{equation*}
                \begin{aligned}
                    C(U_2) &= \frac{C_{0,\text{лин}}}{\sqrt[3]{1 - \frac{U_2}{\varphi_k}}} = \\
                    &= \directlua{
                        tex.sprint("\\frac{"
                                   .. include.value_with_mantice(data.C_0_lin, 2, 0, 0) .. 
                                   "}{\\sqrt[3]{1 - \\frac{"
                                   .. include.value_with_mantice(data.U_2, 2, 0, 0) .. 
                                   "}{"
                                   .. include.value_with_mantice(data.phi_k_lin, 2, 0, 0) .. 
                                   "}}} = "
                                   .. include.value_with_mantice(data.C_lin_U2, 2, 0, 0))
                    } \text{ Ф}
                \end{aligned}
            \end{equation*}

            Ответ: $C(-5\text{ В}) = \directlua{
                tex.sprint(include.value_with_mantice(data.C_lin_U1, 2, 0, 0))
            } \text{ Ф}$, $C(+0.5\text{ В}) = \directlua{
                tex.sprint(include.value_with_mantice(data.C_lin_U2, 2, 0, 0))
            } \text{ Ф}$.
    \end{minipage}\end{minipage}
\end{figure}

\end{document}
